\section{Introdução}

A acessibilidade urbana é tradicionalmente avaliada por métricas que privilegiam a eficiência das redes de transporte, muitas vezes desconsiderando como desigualdades socioespaciais moldam o acesso a serviços cotidianos. Como demonstram \cite{FolGallez2014}, essa abordagem falha em capturar barreiras enfrentadas por populações vulneráveis, especialmente em cidades marcadas por segregação espacial. As autoras argumentam que análises de acessibilidade devem integrar não apenas infraestrutura, mas também a distribuição de serviços essenciais e características da população.

Este artigo propõe uma análise integrada, combinando um modelo de regressão (Experimento 1) com clusterização socioespacial (Experimento 2) para identificar perfis de desigualdade no acesso a pé em Curitiba. Enquanto o Experimento 1 quantifica o impacto de variáveis socioeconômicas nos tempos de acesso, o Experimento 2 identifica padrões espaciais, revelando territórios onde políticas públicas devem ser priorizadas. Além de integrar abordagens metodológicas complementares, este estudo oferece evidências sobre justiça espacial ao demonstrar como desigualdades socioeconômicas moldam a acessibilidade temporal a pé, revelando territórios prioritários para intervenção urbana e oferecendo um modelo que combina técnicas quantitativas e espaciais de forma replicável.

O artigo está organizado em quatro seções principais. Na Seção \ref{sec:related_work}, estudos prévios sobre acessibilidade urbana e análises socioespaciais são revisados. A Seção \ref{sec:methodology} detalha os dois experimentos conduzidos: (1) análise de regressão para identificar fatores socioeconômicos associados ao tempo de acesso e (2) clusterização baseada em PCA para mapear padrões espaciais. A Seção \ref{sec:results}, apresenta os coeficientes significativos do modelo de regressão e uma caracterização dos clusters identificados. Por fim, contribuições, limitações e direções futuras são discutidas na Seção \ref{sec:considerations}.

